\begin{requirementspage}{op}{op integration}{
    Keyao Shen, Jeb Bearer, Verity Coltman
}
    \begin{requirement}{General Rollup Integration}
        OP integrations should meet all high-level requirements of a general
        \page{rollup integration}, with the exception that permissionless batching can be supported
        in a later version. Below are the specific requirements for each component.

        \acceptance All the acceptance criteria of those requirements. If after benchmarks, some
        time requirements don't seem possible, it's okay to update the specific numbers, but the
        general requirements shouldn't be compromised.

        \acceptance
            \begin{itemize}
                \item Benchmark the time a transaction appears in Espresso. The maximum load should
                be one non-empty block every 2s, which applies to the following benchmarks as well.
                \item Benchmark the time a transaction is confirmed on Espresso.
                \item Benchmark the time the batcher is ready to submit a transaction to L1 after
                Espresso confirms it.
                \item Benchmark the time a transaction is confirmed on L1.
            \end{itemize}
    \end{requirement}

    \begin{requirement}{Batch Authentication}
        The authenticator contract must correctly authenticate batches only when they are being
        posted from within a \tee using the correct code and configuration.

        \acceptance Try to register a batcher key in the authenticator contract with an invalid TEE
        attestation. Check that the contract reverts.

        \acceptance Try to authenticate a batch using a signing key that is not registered.
        Check that the contract reverts.
    \end{requirement}

    \begin{requirement}{Batch Posting}
        The batcher must authenticate itself, e.g., by a signature on transactions, before
        submitting batches to L1.

        \acceptance Try to submit a batch to the inbox before authenticating it with the
        authenticator contract. Check that the inbox reverts. Check that the batch is not processed
        by the derivation pipeline.

        \acceptance Run the batcher inside a correctly configured enclave. Check that it
        successfully sends batches to the inbox without reverts.
    \end{requirement}

    \begin{requirement}{Stateless Batcher}
        Operations should be stateless to prevent attacker manipulation that may lead to duplicated
        or mismatched batches. In our threat model, we assume that an attacker may have the capacity
        to manipulate persistent local storage but not the process memory. This means that if our
        operations rely on persistent local storage, an attacker could potentially alter or corrupt
        the data. A stateless approach minimizes the risk of such manipulations.

        \acceptance Run the rollup and randomly restart the batcher. Check the liveness of the
        rollup, and the consistency of Espresso confirmations and L1 confirmations. We don't need
        to clear persistent storage because the original Optimism code isn't and our integration
        work shouldn't use any.
    \end{requirement}

    \begin{requirement}{Derivation From Finalized L1} \label{derivation-from-finalized-l1}
        The caffeinated node should derive from finalized L1 data only, so it is not affected by L1
        reorgs.

        \acceptance Post batches to both Espresso and the L1 then force the L1 to reorg back to an
        earlier state in which those batches have not been posted. Wait for some time and check
        that the batches are eventually posted to the L1 again, in the same order as they were
        originally sequenced, as if the reorg did not happen.

        \acceptance Run a caffeinated node while the L1 is undergoing reorgs, the caffeinated node
        itself never surfaces reorgs or inconsistent state. This should hold even when the
        sequencer is (maliciously) sequencing OP blocks whose L1 origin points to unfinalized L1
        blocks.

        \acceptance Simulate a sequencer that sends a batch derived from an L1 block that hasn't
        been finalized yet. The batcher should not submit it to the L1 before the finalization,
        but later when the L1 eventually finalizes the block, the batcher should include the batch.

        \acceptance Simulate a sequencer that sends a batch derived from an L1 block that hasn't
        been finalized and an L1 reorg that causes a different block being finalized. The batcher
        should ignore the batch derived from the first block.

        Note: The last three acceptance criteria above are for chains like Celo that have a
        a configuration that allows the OP node to only fetch blocks derived from finalized L1
        states. See \subsect{\currentChapter:Additional Considerations:L1 Reorg Handling} for more
        considerations.

    \end{requirement}

    \begin{requirement}{Pipeline Enhancement}
        The derivation pipeline should fully prevent unsuccessful batches from being included on L2.

        \acceptance Attempt to post a batch that fails in the batch inbox contract. The revert
        transactions should not be included by the derivation pipeline.
    \end{requirement}

\end{requirementspage}
